\documentclass[11pt,a4paper]{article}
\usepackage[T1]{fontenc}
\usepackage[utf8]{inputenc}
\usepackage[english]{babel}
\usepackage{amsmath,amssymb}
\usepackage[margin=2cm]{geometry}
\setlength{\parindent}{0pt}
\setlength{\emergencystretch}{3em}
\allowdisplaybreaks

\title{About \texttt{maps\_by\_cases.tex}}
\author{}
\date{}
\begin{document}
\maketitle
This note explains what the test data in \texttt{maps\_by\_cases.tex} is used for. That file is input for the engine in \texttt{main.c}; it is read by the test \texttt{test\_set\_theory.c} in this folder.

\begin{enumerate}
\item \textbf{Purpose.} The file holds 33 true statements about maps between number sets defined by cases: on the remainder in a division (by 2 and by 3), on an order ($n \ge 1$, $n \le 2$, $x \ge 0$), on a point ($x = 1$), on being an integer ($x \in \mathbb{Z}$), on a parameter ($k \mid n$, with $k \in \mathbb{N}$ declared among the hypotheses), and with absolute values ($|n - 2|$, $2x + |x|$), which are read as definitions by cases; items 13--17 have values in pairs, $\mathbb{Z} \to \mathbb{Z}^{2}$ or $\mathbb{Z} \to \mathbb{Z} \times \mathbb{N}_{0}$, and one compares $f(x)$ with pairs written out; items 18 and 19 compare a variable over $\mathbb{Z} \times \mathbb{N}_{0}$ or $\mathbb{N} \times \mathbb{Z}$ with a pair written out, where the two sides differ, so that they become false if the reader exchanges the two components (their false twins are items 22 and 23 of \texttt{false\_maps\_by\_cases.tex}); items 20--22 claim a closed form for the iterates $f^{n}(x)$, items 23--26 are maps from $\mathbb{R}$ (one by cases, one with an absolute value), and items 27--30 have parameters, declared one at a time or together ($a, b \in \mathbb{Q}$); in item 29 the bound $k$ of the conclusion hides the parameter $k$. The last three are maps from $\mathbb{C}$ given by $f(a + ib)$; the last one compares $f(z)$ with $z$ itself, so that $f$ must receive the two components of $z$ in their order. Each claims, but for items 20--22, injectivity, surjectivity or their negations, and must be proved. They are not copies of the exercises of \texttt{maps.tex}: they are other facts of the same kinds, so that the rules are tested on statements the books do not contain.

\item \textbf{What is being tested.} The reading of definitions by cases and of parameters (paragraph 17s of \texttt{parser.c}) and the rules 18v of \texttt{main.c}: that the cases define a map (they cover the domain, overlapping cases agree, the values lie in the codomain), that the case of each $f(t)$ is proved before it is used, that a variable is divided into the cases the proof needs (a remainder, a side of a number, a point, being an integer or a multiple), and that a preimage is looked for case by case. For the maps with values in pairs, the reading of a pair as two numbers (paragraph 17s): an equality between pairs is the equality of the components, and a variable over a set of pairs is two variables, its components (paragraph 17t-2); and the rule 18v-7, by which two conditions that cannot fail together give $1 = 0$ as a combination. For the iterates, their reading (17s) and paragraph 26c of \texttt{main.c}: for a fixed $x$ the numbers $f^{n}(x)$ are a sequence defined by recursion, $s_{1} = f(x)$ and $s_{n + 1} = f(s_{n})$, and the closed form is proved for it as in chapter 26. For the maps from $\mathbb{R}$, rule 18v-8: each $f(t)$ is replaced by its cases, and the statements without the map, in one or two real letters, are decided by the sign table (18p, 18t), after it has proved that the cases cover the domain, agree where they overlap and give values in the codomain. For the parameters: the bound letter that hides a parameter is renamed (17s), the proof divides the values of a parameter into the cases it needs (18v-2: $a = 0$ or not), a witness can be a parameter plus a number (18z: $y = b + 1$, with $b$ in the rule of $f$), and $a x = 1$ forces $a$ to divide $1$ (18v-10). For $\mathbb{C}$: a variable over $\mathbb{C}$ is a pair of real letters, the real part and the coefficient of $i$ (17t-2), the rule is split into real and imaginary part with $i^{2} = -1$ (17s), the definition is unfolded as for $\mathbb{R}$, and an ``exists among the reals at the head of a claim is given rational values, after which the sign table decides what is left (18v-11). Every written step is checked independently, and the statement is checked again on numbers (chapter 19).

\item \textbf{How to run it.} From the top folder of the repository:
\begin{verbatim}
cc -std=c11 -O2 -Wall -Wextra -o /tmp/test_set_theory \
    tests/set_theory/test_set_theory.c && /tmp/test_set_theory
\end{verbatim}

\item \textbf{Expected result.} One \texttt{PASS} line for the file: every statement is proved, with its written proof. The false statements of the same kinds are in \texttt{false\_maps\_by\_cases.tex}.
\end{enumerate}
\end{document}
